\section{Soustraire des nombres relatifs}

\begin{propbox}
\bb{Soustraire un nombre, c'est ajouter son opposé.}

\[
a-b=a+(-b)
\]
\end{propbox}

\begin{exemplebox}
\[
(+7)-(+3)
= (+7)+(-3)
= +4
\]

\[
(+7)-(-3)
= (+7)+(+3)
= +10
\]

\[
(-5)-(+4)
= (-5)+(-4)
= -9
\]
\end{exemplebox}

\begin{methodebox}
Pour effectuer une soustraction :

\begin{enumerate}
    \item Je transforme la soustraction en addition.
    \item Je prends l'opposé du deuxième nombre.
    \item J'applique la règle de l'addition.
\end{enumerate}
\end{methodebox}

\begin{attentionbox}
Attention aux deux signes !

\[
5-(-3) \neq 5-3
\]

En effet :
\[
5-(-3)=5+3=8.
\]
\end{attentionbox}

